\bookchapter{Kapitel 6: Beschütze dein digitales Leben – Sicherheit und Datenschutz}

\chapterintro{Smartphonesicherheit = Datenschutz + Vermögenssicherheit. In einer Zeit, in der Daten ein Vermögenswert sind, ist der Schutz des Smartphones der Schutz deiner selbst.}

\booksection{Eins. Sperrbildschirm: Die erste physische Verteidigungslinie}

Der Sperrbildschirm ist die erste Verteidigungslinie deines Smartphones – und wird am leichtesten übersehen.

\booksubsection{Sicherheitsstufen-Rangfolge:}

\begin{enumerate}
  \item \textbf{Passwort/PIN}: Höchste Sicherheit, die ultimative Absicherung (wenn Fingerabdruck/Gesicht vergessen).
  \item \textbf{Fingerabdruck}: Perfekte Balance aus Sicherheit und Komfort.
  \item \textbf{Gesichtserkennung}: Am bequemsten, aber die Sicherheit hängt von der Technik ab (3D-Strukturlicht sicher, 2D-Foto leicht auszutricksen).
  \item \textbf{Muster-Entsperrung}: Praktisch, aber es bleiben Spuren zurück, die erraten werden können.
\end{enumerate}

\booksubsection{Starken Sperrbildschirm einrichten:}

\begin{itemize}
  \item Unbedingt ein komplexes Passwort festlegen (nicht Geburtsdatum oder 123456).
  \item Biometrische Entsperrung aktivieren.
  \item Sperrzeit auf „Sofort" oder kürzeste Zeitspanne einstellen.
\end{itemize}

\booksection{Zwei. „Mein Smartphone finden“: Das ultimative Diebstahlschutz-Werkzeug}

\textbf{Muss aktiviert sein!}Es kann:

\begin{itemize}
  \item \textbf{Orten}: Wo das Smartphone hingeblieben ist, auf der Karte finden.
  \item \textbf{Ton abspielen}: Das Smartphone klingeln lassen – ideal zum Suchen zu Hause.
  \item \textbf{Verloren-Modus}: Das Smartphone aus der Ferne sperren und eine Kontaktnummer anzeigen.
  \item \textbf{Fernlöschung}: Die letzte Option – alle Daten löschen (danach ist keine Ortung mehr möglich).
\end{itemize}

Unter [Einstellungen] bei Apple „Mein iPhone suchen“, bei Android „Mein Gerät finden” oder einen markenähnlichen Dienst suchen und sicherstellen, dass er aktiviert (grün) ist.

\booksection{Drei. System-Updates: Patches einspielen, Lücken stopfen}

Der Hauptzweck von System-Updates ist die Behebung von Sicherheitslücken – wie das Flicken von Löchern in einem Haus. Wenn du also eine Update-Aufforderung siehst, aktualisiere möglichst umgehend.

Natürlich können große Versions-Updates (z. B. von iOS 17 auf iOS 18) Kompatibilitätsprobleme mit sich bringen; hier kannst du einige Tage abwarten, bevor du aktualisierst.

\booksection{Vier. Telekommunikationsbetrug erkennen und verhindern}

In einer Zeit rasant fortschreitender KI-Technik ist die traditionelle Überzeugung „Sehen heißt Glauben“ vollständig hinfällig. Betrüger nutzen Deepfake- und andere KI-Techniken, um vom „Gießkannen-Betrug” zur „präzisen Jagd“ überzugehen.

\booksubsection{Die häufigsten Betrugsmaschen im KI-Zeitalter}

\textbf{(1) KI-Tausch von Gesicht und Stimme (Deepfake-Betrug)}\par

Täter sammeln illegal hochauflösende Fotos von Zielpersonen aus sozialen Medien oder nur wenige Sekunden Sprachprobe und klonen mit KI-Deepfake-Technik ein täuschend echtes „Familien-Video“ oder eine „Stimme“. Typische Drehbücher:\textbf{plötzlicher Unfall, dringendes Darlehen, Entführung}und Ähnliches. Solche Betrügereien nutzen die emotionale Bindung von Familie und Freundschaft und sind\textbf{extrem täuschungsstark}.

\textbf{(2) Falscher Amtsträger / falscher Kundenservice}\par

KI kann Stimme und Tonfall von offiziellen Kundendiensten, Polizei und Staatsanwaltschaft präzise imitieren und sogar hochgradig echt wirkende gefälschte Überweisungs-Screenshots, elektronische Amtsdokumente oder Fahndungsschreiben erzeugen. Über\textbf{Panikmache oder Autoritätsdruck}wird das Opfer zur Überweisung oder zur Bildschirmfreigabe verleitet – und die Kontodaten sind weg.

\textbf{(3) KI-Anlagebetrug}\par

Mit KI-Algorithmen werden „garantierter Gewinn“-Screenshots, Interfaces angeblicher Quanten-Auswahlsoftware oder „Meister“-Analysen gefälscht. Mit derlockenden hohen Rendite werden teure Anlagedienste verkauft oder das Opfer bringt sein Geld auf ein Konto der Täter –\textbf{am Ende verschwinden sie mit dem Geld}.

\booksubsection{Der Ausweg – der „Zauberspiegel“-Schutzansatz}

KI kann Stimme und Gesicht klonen, ist aber nicht unfehlbar. Es gibt Schwächen zum Ausnutzen:

\begin{itemize}
  \item \textbf{Bewegungsprüfung}: Fordere beim Videoanruf aktiv einen vorab vereinbarten, zusammenhängenden Bewegungsablauf – etwa „wische schnell mit der Hand übers Gesicht“, „zeig im Profil die Ohrenkontur“ oder „blinzel 5-mal hintereinander schnell“. Heutige Face-Swap-Technik zeigt bei komplexen Verdeckungen, schnellen Bewegungen und Lichtwechsel\textbf{leicht Ruckler, verzerrte Gesichter oder Kantenfehler}.
  \item \textbf{Geheimcode-Prüfung}: Vereinbare und nutze private Prüf-Fragen, die nur ihr beide kennt – etwa „was haben wir beim letzten gemeinsamen Essen bestellt?“ oder „was war das erste Geschenk, das ich dir gemacht habe?“. Zögert die Gegenseite auffällig oder antwortet am Thema vorbei, dann\textbf{sofort das Gespräch beenden und den Sicherheitsplan starten}.
  \item \textbf{Zweifach-Nachfrage}: Folge strikt dem Grundsatz der Cyberseinheit der öffentlichen Sicherheit,\textbf{„nicht-glauben, nicht-überweisen, nicht-preisgeben, nicht-weiterleiten – aber eins verifizieren“}– fremden Nachrichten nicht trauen, nicht einfach überweisen, keine sensiblen Daten preisgeben, keine verdächtigen Inhalte weiterleiten. Bei allem, was mit Geld zu tun hat, unbedingt über die\textbf{unbedingt über die bereits bekannte Originalnummer zurückruft oder die Person persönlich bestätigt}. Bei Überweisungen wird der Modus „24-Stunden-Ziel" empfohlen, um ausreichend Zeit für Betrugserkennung und Rückerstattung zu haben.
\end{itemize}

\booksubsection{Klassische Maschen (weiterhin wachsam bleiben):}

\begin{itemize}
  \item \textbf{Gewinn-/Rebate-Lockrufe}: Lockmittel – erst ein kleiner Sweetener, dann der große Betrug.
  \item \textbf{Romance-Scamming (Betrug über Online-Beziehungen)}: „Schweine-Schlachten“-Schema – erst Gefühle aufbauen, dann Geld abziehen.
  \item \textbf{Ausgabe als Bekannter}: Geklauter Account, dann Geld leihen (im KI-Zeitalter noch schwerer zu erkennen).
\end{itemize}

\booksubsection{Der eingebaute Schutz des Handys:}

\begin{itemize}
  \item Belästigungsanrufe blockieren.
  \item Schutz vor SMS-Code-Diebstahl.
  \item Sicherheitsprüfung von SMS-Links.
\end{itemize}

\booksubsection{Der Anti-Betrugs-Merkvers:}

\textbf{Nicht leichtgläubig sein, nichts verraten, nichts antippen, nichts überweisen – bei Zweifel erst verifizieren!}Bei einer „Notlage“ erst 5 Minuten ruhig bleiben – fast jeder Betrug fliegt in genau diesem Moment auf.

\booksection{Schutz vor „Doxxing“ und Social-Engineering: die Privatkette kappen}

Der Kern des „Doxxing“ (Social-Engineering-Angriffs) liegt darin,\textbf{systematisch digitale Splitter einer Person zu sammeln, zu analysieren und zu einem vollständigen Persönlichkeitsbild zusammenzusetzen}. Entscheidend für den Schutz: klare rote Linien für die eigene Privatsphäre und ein mehrschichtiges Schutzsystem.

\booksubsection{Bilder anonymisieren}

Allgemeine Werkzeuge: ExifTool, Metapho usw. können EXIF anzeigen/bearbeiten.\textbf{Metadaten von Fotos (EXIF) vollständig löschen}– Aufnahmezeitpunkt, Gerätemodell, GPS-Koordinaten und andere sensible Daten. Und die\textbf{Hausnummern, Kfz-Kennzeichen, Barcode der Fahrkarte, Ausweisnummer}auf den Bildern konsequent verpixeln oder unscharf machen, damit keine Bewegungsprofile durchsickern.

\begin{itemize}
  \item iOS: [Einstellungen] > [Datenschutz \& Sicherheit] > [Kamera] – dort „Genauer Standort“ deaktivieren; beim Teilen „Nur aktuelles Foto“ wählen, das streift die Metadaten ab.
  \item Android: Einige Album-Apps (z. B. Google Fotos) fragen beim Teilen, ob Standortinformationen beibehalten werden sollen.
  \item Allgemein: ExifTool, Metapho und Ähnliches zeigen und bearbeiten EXIF-Daten.
\end{itemize}

\booksubsection{App-Berechtigungen zusammenstreifen}

Bei der Anfrage nach „Kontaktberechtigung“ erst fragen: Muss diese App wirklich alle meine Kontakte kennen?\textbf{Mikrofon, Kontakte, Standort}und andere unnötige Berechtigungen entziehen – sonst „lauscht“ die App im Hintergrund auf Alltagsdialoge oder sammelt dein Beziehungsnetz. Apps, die lange ungenutzt hohe Rechte halten, promptly deinstallieren – das mindert die Risikopunkte an der Quelle.

\begin{itemize}
  \item Bei der Anfrage „Kontakte-Berechtigung“ erst fragen: Muss diese App wirklich alle meine Kontakte kennen?
  \item Apps, die du drei Monate nicht genutzt hast: entschlossen deinstallieren. Berechtigungen werden mit der Zeit enger gezogen, nicht lockerer.
\end{itemize}

\booksubsection{Smarte Geräte absichern}

IoT-Geräte wie Smartkameras, Smartlocks und Smartlautsprecher müssen beim ersten Gebrauch\textbf{sofort das Standardpasswort ändern}– ein starkes, eigenständiges Passwort mit Groß-/Kleinbuchstaben, Ziffern und Sonderzeichen setzen; selten genutzte Fernzugriffsfunktionen ausgeschaltet lassen, Firmware-Updates regelmäßig prüfen und Sicherheitslücken promptly schließen,\textbf{damit Hacker nicht ins familiäre Privatleben spähen}.

\begin{tcbwarn}{Warnung aus der Praxis}
2022 meldete die Polizei eines Bezirks, dass Hacker durch das Knacken des Standardpassworts einer Haushaltskamera-Marke (viele Nutzer hatten es nie geändert) Tausende von Schlafzimmerszenen gestohlen und online verkauft hatten.\textbf{Standardpasswort nicht ändern = die Haustür sperrangelweit offen}.

\end{tcbwarn}

\booksection{Privatsphäre: deine Daten sind dein Eigentum}

Fotos, Chatverläufe, Bewegungsprofile im Handy sind im Kern dein „digitales Vermögen“. Einmal ausgeleckt: im besten Fall Werbebelästigung, im schlimmsten Fall Doxxing, Betrug bis hin zu Erpressung.

\booksubsection{: Ab iOS 14.5+ muss eine App vor der Verwendung von IDFA fragen; Android kann die „Werbe-ID“ deaktivieren.}

\begin{itemize}
  \item \textbf{Standort}: Für Taxi/Karten „Immer erlauben“; für andere Apps „Nur bei Verwendung“.
  \item \textbf{Kontakte}: Nur wenn nötig (etwa WeChat-Freunde hinzufügen) – sonst grundsätzlich verweigern.
  \item \textbf{Galerie / alle Fotos}: Ab iOS 14+ können „Ausgewählte Fotos" freigegeben werden; Android 13+ unterstützt ebenfalls feingranulare Berechtigungen.
  \item \textbf{Mikrofon / Kamera}.
  \item \textbf{Hintergrund-Tracking}: Ab iOS 14.5 müssen Apps vor der Nutzung der IDFA um Erlaubnis fragen; bei Android lässt sich die „Werbe-ID“ abschalten.
\end{itemize}

\booksubsection{Datenschutz-Praxis}

\begin{itemize}
  \item \textbf{.}: [Einstellungen] > [Datenschutz \& Sicherheit] > [App-Datenschutz-Bericht], zeigt welche Apps heimlich auf Standort, Kamera oder Mikrofon zugreifen.
  \item \textbf{Berechtigungsverlauf (Android)}: [Einstellungen] > [Datenschutz] > [Berechtigungsverlauf], zeigt wer wann sensible Berechtigungen aufgerufen hat.
  \item \textbf{Werbetracking einschränken}: In den Systemeinstellungen „Werbe-ID" deaktivieren / „Werbetracking einschränken", um „präzise Werbebelästigung" zu reduzieren.
\end{itemize}

\booksection{Sieben. „Erste Hilfe“-Checkliste nach Verlust oder Diebstahl}

Wenn das Smartphone verloren ist, nicht in Panik verfallen – in dieser Reihenfolge vorgehen:

\begin{enumerate}
  \item \textbf{SIM-Karte sofort sperren}: Netzbetreiber kontaktieren und Verlust melden, um Missbrauch für Bestätigungscodes zu verhindern.
  \item \textbf{Aus der Ferne sperren}: Über Cloud-Dienst (iCloud / Marken-Cloud) den „Verloren-Modus" aktivieren.
  \item \textbf{Wichtige Passwörter ändern}: Alle wichtigen Konten, die auf dem Smartphone eingeloggt waren (Zahlung, Bank, E-Mail), sofort ändern.
  \item \textbf{Anzeige erstatten}: Zur Polizei gehen und die IMI-Nummer notieren. Falls das Smartphone wiedergefunden wird, ist die Identität leichter nachzuweisen.
  \item \textbf{Familie und Freunde benachrichtigen}: Verhindern, dass Betrüger in deinem Namen Geld fordern.
  \item \textbf{(Zuletzt) Fernlöschung}: Erst durchführen, wenn die Wiederbeschaffung aussichtslos ist; danach ist die Ortung nicht mehr möglich.
\end{enumerate}

\booksection{Acht. SIM-Karte und eSIM-Sicherheit}

\begin{itemize}
  \item \textbf{SIM-PIN einrichten}: Nach der Aktivierung ist bei jedem Neustart oder Kartenwechsel die PIN-Eingabe erforderlich, um das Umgehen der Sperre durch „Karte entnehmen und in anderes Gerät einlegen" zu verhindern.
  \item \textbf{eSIM-Vorteile}: Keine physische Karte entfernbar, höhere Sicherheit; unterstützt Mehrfachnummern-Umschaltung, sehr praktisch auf Reisen.
  \item \textbf{Risiko der SIM-Karten-Zweitbezug-Betrug}: Vor der Netzbetreiber-Masche „Zweitbezug" warnen – Betrüger legen mit deiner Identität eine neue SIM-Karte zu und übernehmen so deine Handynummer. Die Einstellung „Kartenersatz erfordert persönliches Erscheinen des Inhabers mit Ausweis in der Filiale" ist die wirksamste Verteidigungslinie.
\end{itemize}

\begin{tcbwarn}{Was du niemals tun solltest}
\begin{itemize}
  \item Schwache Passwörter wie Geburtsdatum / 123456 / qwerty verwenden;
  \item Bestätigungscodes an irgendjemanden weitergeben (auch nicht an „Kundendienst");
  \item Auf unbekannte Links in SMS klicken, besonders bei ausländischen Nummern mit 00-Beginn;
  \item Unter öffentlichem WLAN Online-Banking oder Zahlungen nutzen;
  \item Fremden „Bildschirmfreigabe" oder „Fernsteuerung" gewähren.
\end{itemize}

\end{tcbwarn}

\begin{tcbpractice}{Versuche, „Sperren aus der Ferne" zu aktivieren und zu simulieren}
\begin{enumerate}
  \item Stelle sicher, dass die Funktion „Mein Smartphone finden" aktiviert ist.
  \item Öffne auf einem Computer-Browser oder einem anderen Smartphone den entsprechenden Cloud-Dienst (iCloud.com oder die Marken-Cloud-Website).
  \item Finde die Funktion „Gerät finden" und versuche, eine „Ton abspielen"- oder „Sperren"-Anweisung an dein eigenes Smartphone zu senden, um den Ablauf kennenzulernen.
\end{enumerate}

\end{tcbpractice}
